\documentclass[preview]{standalone}

\usepackage{amsmath}
\usepackage{amssymb}
\usepackage{stellar}
\usepackage{definitions}
\usepackage{bettelini}

\begin{document}

\title{Stellar}
\id{linearalgebra-kernel-image}
\genpage

\section{Kernel}

\begin{snippetdefinition}{linear-transformation-kernel-definition}{Kernel of a linear transformation}
    Let \(V\) and \(W\) be \vectorspace[vector spaces] over a \field
    \(\mathbb{F}\), and let
    \(f\in\setoflineartransformations(V,W)\). The \emph{kernel} of \(f\) is
    the \set
    \[
        \ltkernel f
        =
        \{v\in V\suchthat f(v)=0_W\}.
    \]
\end{snippetdefinition}

\begin{snippetexample}{linear-transformation-basic-kernels-example}{Basic kernels}
    Let \(V\) and \(W\) be \vectorspace[vector spaces] over a \field
    \(\mathbb{F}\). If
    \[
        0\colon V\fromto W
    \]
    is the zero \lineartransformation, then
    \[
        \ltkernel 0=V.
    \]
    If
    \[
        \identityfunc_V\colon V\fromto V
    \]
    is the identity \lineartransformation, then
    \[
        \ltkernel \identityfunc_V=\{0_V\}.
    \]
\end{snippetexample}

\begin{snippetexample}{complex-linear-map-kernel-example}{A kernel in \(\mathbb{C}^3\)}
    Consider the \lineartransformation
    \[
        f\colon \complexnumbers^3\fromto\complexnumbers,
        \qquad
        f(z_1,z_2,z_3)=z_1+\mathrm{i}z_2+2\mathrm{i}z_3.
    \]
    Then
    \[
        \ltkernel f
        =
        \{(z_1,z_2,z_3)\in\complexnumbers^3
        \suchthat
        z_1+\mathrm{i}z_2+2\mathrm{i}z_3=0\}.
    \]
    The equation gives
    \[
        z_1=-\mathrm{i}z_2-2\mathrm{i}z_3,
    \]
    so
    \[
        \ltkernel f
        =
        \linearspan\bigl((-\mathrm{i},1,0),(-2\mathrm{i},0,1)\bigr).
    \]
\end{snippetexample}

\begin{snippetexample}{polynomial-derivative-kernel-example}{The kernel of the derivative}
    Let
    \[
        D\colon \realnumbers[x]\fromto\realnumbers[x],
        \qquad
        D(p)=p',
    \]
    be the derivative \lineartransformation. Then
    \[
        \ltkernel D
        =
        \{p\in\realnumbers[x]\suchthat p'=0\}
    \]
    is the \set of constant \polynomial[polynomials].
\end{snippetexample}

\begin{snippetexample}{backward-shift-kernel-example}{Kernel of the backward shift}
    Let \(V=\mathbb{F}^{\naturalnumbers}\), and let
    \[
        b_s\colon V\fromto V,
        \qquad
        b_s(x_1,x_2,x_3,\ldots)=(x_2,x_3,x_4,\ldots).
    \]
    Then
    \[
        \ltkernel b_s
        =
        \{(x_1,0,0,\ldots)\suchthat x_1\in\mathbb{F}\}
        =
        \linearspan\bigl((1,0,0,\ldots)\bigr).
    \]
\end{snippetexample}

\begin{snippetproposition}{kernel-of-linear-transformation-is-subspace}{The kernel is a subspace}
    Let \(V\) and \(W\) be \vectorspace[vector spaces] over a \field
    \(\mathbb{F}\), and let
    \(f\in\setoflineartransformations(V,W)\). Then \(\ltkernel f\) is a
    linear subspace of \(V\).
\end{snippetproposition}

\begin{snippetproof}{kernel-of-linear-transformation-is-subspace-proof}{kernel-of-linear-transformation-is-subspace}{The kernel is a subspace}
    Since \(f(0_V)=0_W\), we have \(0_V\in\ltkernel f\).

    If \(u,v\in\ltkernel f\), then \(f(u)=0_W\) and \(f(v)=0_W\). By
    linearity,
    \[
        f(u+v)=f(u)+f(v)=0_W+0_W=0_W,
    \]
    so \(u+v\in\ltkernel f\).

    If \(\lambda\in\mathbb{F}\) and \(v\in\ltkernel f\), then
    \[
        f(\lambda v)=\lambda f(v)=\lambda 0_W=0_W,
    \]
    so \(\lambda v\in\ltkernel f\). By the linear subspace criterion,
    \(\ltkernel f\) is a linear subspace of \(V\).
\end{snippetproof}

\begin{snippetproposition}{linear-transformation-injective-kernel-criterion}{Injectivity and kernel}
    Let \(V\) and \(W\) be \vectorspace[vector spaces] over a \field
    \(\mathbb{F}\), and let
    \(f\in\setoflineartransformations(V,W)\). Then \(f\) is
    \injective if and only if
    \[
        \ltkernel f=\{0_V\}.
    \]
\end{snippetproposition}

\begin{snippetproof}{linear-transformation-injective-kernel-criterion-proof}{linear-transformation-injective-kernel-criterion}{Injectivity and kernel}
    \begin{iffproof}
        \item Suppose that \(f\) is \injective. Since \(f(0_V)=0_W\), the
        zero \vector belongs to \(\ltkernel f\). If \(v\in\ltkernel f\), then
        \(f(v)=0_W=f(0_V)\). Injectivity gives \(v=0_V\). Hence
        \(\ltkernel f=\{0_V\}\).

        \item Suppose that \(\ltkernel f=\{0_V\}\). Let \(u,v\in V\) and
        assume \(f(u)=f(v)\). Then
        \[
            f(u-v)=f(u)-f(v)=0_W,
        \]
        so \(u-v\in\ltkernel f\). Hence \(u-v=0_V\), and therefore
        \(u=v\). Thus \(f\) is \injective.
    \end{iffproof}
\end{snippetproof}

\section{Image}

\begin{snippetdefinition}{linear-transformation-image-definition}{Image of a linear transformation}
    Let \(V\) and \(W\) be \vectorspace[vector spaces] over a \field
    \(\mathbb{F}\), and let
    \(f\in\setoflineartransformations(V,W)\). The \emph{image} of \(f\) is
    the \set
    \[
        \ltimage f
        =
        \{w\in W\suchthat \exists v\in V,\ f(v)=w\}.
    \]
    Equivalently,
    \[
        \ltimage f=f(V).
    \]
\end{snippetdefinition}

\begin{snippetexample}{linear-transformation-basic-images-example}{Basic images}
    Let \(V\) and \(W\) be \vectorspace[vector spaces] over a \field
    \(\mathbb{F}\). If \(0\colon V\fromto W\) is the zero
    \lineartransformation, then
    \[
        \ltimage 0=\{0_W\}.
    \]
    If \(\identityfunc_V\colon V\fromto V\) is the identity
    \lineartransformation, then
    \[
        \ltimage \identityfunc_V=V.
    \]
\end{snippetexample}

\begin{snippetexample}{r2-r3-linear-map-image-example}{An image in \(\mathbb{R}^3\)}
    Consider the \lineartransformation
    \[
        f\colon \realnumbers^2\fromto\realnumbers^3,
        \qquad
        f(x,y)=(2x,5y,x+y).
    \]
    Since
    \[
        f(x,y)=x(2,0,1)+y(0,5,1),
    \]
    we have
    \[
        \ltimage f
        =
        \linearspan\bigl((2,0,1),(0,5,1)\bigr)
        \subseteq\realnumbers^3.
    \]
\end{snippetexample}

\begin{snippetexample}{derivative-image-example}{The image of the derivative}
    Let
    \[
        D\colon \realnumbers[x]\fromto\realnumbers[x],
        \qquad
        D(p)=p',
    \]
    be the derivative \lineartransformation. Then
    \[
        \ltimage D=\realnumbers[x],
    \]
    because every real \polynomial has a polynomial primitive.
\end{snippetexample}

\begin{snippetproposition}{linear-transformation-image-is-subspace}{The image is a subspace}
    Let \(V\) and \(W\) be \vectorspace[vector spaces] over a \field
    \(\mathbb{F}\), and let
    \(f\in\setoflineartransformations(V,W)\). Then \(\ltimage f\) is a
    linear subspace of \(W\).
\end{snippetproposition}

\begin{snippetproof}{linear-transformation-image-is-subspace-proof}{linear-transformation-image-is-subspace}{The image is a subspace}
    Since \(f(0_V)=0_W\), we have \(0_W\in\ltimage f\).

    Let \(w_1,w_2\in\ltimage f\). Then there exist \(v_1,v_2\in V\) such
    that \(w_1=f(v_1)\) and \(w_2=f(v_2)\). By linearity,
    \[
        w_1+w_2=f(v_1)+f(v_2)=f(v_1+v_2),
    \]
    so \(w_1+w_2\in\ltimage f\).

    If \(\lambda\in\mathbb{F}\) and \(w\in\ltimage f\), choose
    \(v\in V\) with \(w=f(v)\). Then
    \[
        \lambda w=\lambda f(v)=f(\lambda v),
    \]
    so \(\lambda w\in\ltimage f\). By the linear subspace criterion,
    \(\ltimage f\) is a linear subspace of \(W\).
\end{snippetproof}

\begin{snippetproposition}{linear-transformation-surjective-image-criterion}{Surjectivity and image}
    Let \(V\) and \(W\) be \vectorspace[vector spaces] over a \field
    \(\mathbb{F}\), and let
    \(f\in\setoflineartransformations(V,W)\). Then \(f\) is
    \surjective if and only if
    \[
        \ltimage f=W.
    \]
\end{snippetproposition}

\begin{snippetproof}{linear-transformation-surjective-image-criterion-proof}{linear-transformation-surjective-image-criterion}{Surjectivity and image}
    The \function \(f\) is \surjective precisely when every \(w\in W\) is
    equal to \(f(v)\) for at least one \(v\in V\). This is exactly the
    assertion that every \(w\in W\) belongs to \(\ltimage f\), namely
    \(\ltimage f=W\).
\end{snippetproof}

\section{Invertibility}

\begin{snippetdefinition}{invertible-linear-transformation-definition}{Invertible linear transformation}
    Let \(V\) and \(W\) be \vectorspace[vector spaces] over a \field
    \(\mathbb{F}\). A \lineartransformation
    \(f\in\setoflineartransformations(V,W)\) is \emph{invertible} if there
    exists \(g\in\setoflineartransformations(W,V)\) such that
    \[
        g\circ f=\identityfunc_V
        \qquad
        \text{and}
        \qquad
        f\circ g=\identityfunc_W.
    \]
    In that case \(g\) is called the inverse of \(f\).
\end{snippetdefinition}

\begin{snippetproposition}{linear-transformation-inverse-uniqueness}{Uniqueness of the inverse}
    Let \(V\) and \(W\) be \vectorspace[vector spaces] over a \field
    \(\mathbb{F}\). If
    \(f\in\setoflineartransformations(V,W)\) is invertible, then its inverse
    is unique.
\end{snippetproposition}

\begin{snippetproof}{linear-transformation-inverse-uniqueness-proof}{linear-transformation-inverse-uniqueness}{Uniqueness of the inverse}
    Let \(g_1,g_2\in\setoflineartransformations(W,V)\) be inverses of \(f\).
    Then
    \[
        g_1\circ f=\identityfunc_V,
        \qquad
        f\circ g_2=\identityfunc_W.
    \]
    Hence
    \[
        g_1
        =
        g_1\circ\identityfunc_W
        =
        g_1\circ(f\circ g_2)
        =
        (g_1\circ f)\circ g_2
        =
        \identityfunc_V\circ g_2
        =
        g_2.
    \]
\end{snippetproof}

\begin{snippettheorem}{linear-transformation-invertible-iff-bijective}{Invertibility criterion}
    Let \(V\) and \(W\) be \vectorspace[vector spaces] over a \field
    \(\mathbb{F}\), and let
    \(f\in\setoflineartransformations(V,W)\). Then \(f\) is invertible if and
    only if \(f\) is both \injective and \surjective.
\end{snippettheorem}

\begin{snippetproof}{linear-transformation-invertible-iff-bijective-proof}{linear-transformation-invertible-iff-bijective}{Invertibility criterion}
    \begin{iffproof}
        \item Suppose that \(f\) is invertible, and let \(f^\inversefunction\)
        be its inverse. If \(f(v_1)=f(v_2)\), then applying
        \(f^\inversefunction\) gives \(v_1=v_2\), so \(f\) is \injective.
        If \(w\in W\), then
        \[
            w=(f\circ f^\inversefunction)(w)=f(f^\inversefunction(w)),
        \]
        so \(w\in\ltimage f\). Hence \(f\) is \surjective.

        \item Suppose that \(f\) is \injective and \surjective. For each
        \(w\in W\), there is a unique \(v\in V\) such that \(f(v)=w\). Define
        \(g\colon W\fromto V\) by assigning to \(w\) this unique \(v\).
        Then \(f\circ g=\identityfunc_W\) and \(g\circ f=\identityfunc_V\).

        It remains to show that \(g\) is linear. Let \(w_1,w_2\in W\). Since
        \[
            f(g(w_1+w_2))=w_1+w_2
        \]
        and
        \[
            f(g(w_1)+g(w_2))=f(g(w_1))+f(g(w_2))=w_1+w_2,
        \]
        injectivity of \(f\) gives
        \[
            g(w_1+w_2)=g(w_1)+g(w_2).
        \]
        Similarly, for \(\lambda\in\mathbb{F}\),
        \[
            f(g(\lambda w))=\lambda w
            =
            f(\lambda g(w)),
        \]
        so \(g(\lambda w)=\lambda g(w)\). Therefore \(g\) is linear, and
        \(f\) is invertible.
    \end{iffproof}
\end{snippetproof}

\begin{snippetexample}{multiplication-by-x-squared-not-invertible-example}{Multiplication by \(x^2\)}
    The \lineartransformation
    \[
        M_{x^2}\colon \realnumbers[x]\fromto\realnumbers[x],
        \qquad
        M_{x^2}(p)=x^2p,
    \]
    is \injective because \(x^2p=0\) implies \(p=0\). It is not
    \surjective because, for instance,
    \[
        x-1\notin\ltimage M_{x^2}.
    \]
    Hence \(M_{x^2}\) is not invertible.
\end{snippetexample}

\begin{snippetexample}{backward-shift-not-invertible-example}{The backward shift is not invertible}
    Let \(V=\mathbb{F}^{\naturalnumbers}\), and let
    \[
        b_s(x_1,x_2,x_3,\ldots)=(x_2,x_3,x_4,\ldots).
    \]
    This \lineartransformation \(b_s\colon V\fromto V\) is \surjective,
    since
    \[
        (y_1,y_2,y_3,\ldots)=b_s(0,y_1,y_2,y_3,\ldots).
    \]
    It is not \injective because
    \[
        \ltkernel b_s=\linearspan\bigl((1,0,0,\ldots)\bigr)\neq\{0\}.
    \]
    Hence \(b_s\) is not invertible.
\end{snippetexample}

\section{Isomorphisms}

\begin{snippetdefinition}{linear-isomorphism-definition}{Linear isomorphism}
    Let \(V\) and \(W\) be \vectorspace[vector spaces] over the same \field
    \(\mathbb{F}\). A \emph{linear isomorphism} from \(V\) to \(W\) is an
    invertible \lineartransformation
    \[
        f\colon V\fromto W.
    \]
    If there exists a \linearisomorphism from \(V\) to \(W\), then \(V\) and
    \(W\) are called isomorphic, and one writes
    \[
        V\simeq W.
    \]
\end{snippetdefinition}

\begin{snippetexample}{basis-coordinate-isomorphism-example}{Coordinates give an isomorphism}
    Let \(V\) be a finite-dimensional \vectorspace over a \field
    \(\mathbb{F}\), and let
    \[
        \mathcal{B}=\{v_1,\ldots,v_n\}
    \]
    be a \basis of \(V\). The coordinate \lineartransformation
    \[
        \Phi_{\mathcal{B}}\colon V\fromto\mathbb{F}^n
    \]
    defined by
    \[
        \Phi_{\mathcal{B}}(a_1v_1+\cdots+a_nv_n)=(a_1,\ldots,a_n)
    \]
    is a \linearisomorphism. It is \injective because only the zero
    \vector has all coordinates equal to \(0\), and it is \surjective because
    each \((a_1,\ldots,a_n)\in\mathbb{F}^n\) is the coordinate vector of
    \(a_1v_1+\cdots+a_nv_n\).
\end{snippetexample}

\begin{snippettheorem}{finite-dimensional-vector-spaces-classification}{Classification of finite-dimensional vector spaces}
    Let \(V\) and \(W\) be finite-dimensional \vectorspace[vector spaces] over
    the same \field \(\mathbb{F}\). Then
    \[
        V\simeq W
        \qquad\Longleftrightarrow\qquad
        \lineardim V=\lineardim W.
    \]
\end{snippettheorem}

\begin{snippetproof}{finite-dimensional-vector-spaces-classification-proof}{finite-dimensional-vector-spaces-classification}{Classification of finite-dimensional vector spaces}
    \begin{iffproof}
        \item Suppose \(f\colon V\fromto W\) is a \linearisomorphism. Since
        \(V\) is finite-dimensional, let
        \[
            \mathcal{B}=\{v_1,\ldots,v_n\}
        \]
        be a \basis of \(V\). We show that
        \[
            f(\mathcal{B})=\{f(v_1),\ldots,f(v_n)\}
        \]
        is a \basis of \(W\).

        If
        \[
            \lambda_1f(v_1)+\cdots+\lambda_nf(v_n)=0_W,
        \]
        then
        \[
            f(\lambda_1v_1+\cdots+\lambda_nv_n)=0_W=f(0_V).
        \]
        Since \(f\) is \injective,
        \[
            \lambda_1v_1+\cdots+\lambda_nv_n=0_V.
        \]
        The \basis \(\mathcal{B}\) is \linearlyindependent, so all
        coefficients vanish. Thus \(f(\mathcal{B})\) is
        \linearlyindependent.

        If \(w\in W\), then \(w=f(v)\) for some \(v\in V\). Writing
        \(v=a_1v_1+\cdots+a_nv_n\), linearity gives
        \[
            w=a_1f(v_1)+\cdots+a_nf(v_n).
        \]
        Thus \(f(\mathcal{B})\) generates \(W\). Therefore
        \(f(\mathcal{B})\) is a \basis of \(W\), so
        \(\lineardim W=n=\lineardim V\).

        \item Suppose \(\lineardim V=\lineardim W=n\). Let
        \[
            \mathcal{B}=\{v_1,\ldots,v_n\},
            \qquad
            \mathcal{C}=\{w_1,\ldots,w_n\}
        \]
        be \basis[bases] of \(V\) and \(W\). By the linear extension theorem,
        there is a unique \lineartransformation \(f\colon V\fromto W\) such
        that \(f(v_i)=w_i\) for every \(i\). Its inverse is the unique
        \lineartransformation \(g\colon W\fromto V\) with \(g(w_i)=v_i\) for
        every \(i\). Therefore \(V\simeq W\).
    \end{iffproof}
\end{snippetproof}

\begin{snippettheorem}{finite-dimensional-injectivity-surjectivity-equivalence}{Injectivity, surjectivity and dimension}
    Let \(V\) and \(W\) be finite-dimensional \vectorspace[vector spaces] over
    a \field \(\mathbb{F}\), and assume
    \[
        \lineardim V=\lineardim W.
    \]
    For \(f\in\setoflineartransformations(V,W)\), the following are
    equivalent:
    \begin{enumerate}
        \item \(f\) is invertible;
        \item \(f\) is \injective;
        \item \(f\) is \surjective.
    \end{enumerate}
\end{snippettheorem}

\begin{snippetproof}{finite-dimensional-injectivity-surjectivity-equivalence-proof}{finite-dimensional-injectivity-surjectivity-equivalence}{Injectivity, surjectivity and dimension}
    An invertible \lineartransformation is \injective and \surjective.
    Conversely, a \lineartransformation that is both \injective and
    \surjective is invertible.

    It remains to show that injectivity and surjectivity imply one another
    under the dimension hypothesis. Put
    \(\lineardim V=\lineardim W=n\), and let
    \[
        \mathcal{B}=\{v_1,\ldots,v_n\}
    \]
    be a \basis of \(V\).

    If \(f\) is \injective, then
    \[
        f(\mathcal{B})=\{f(v_1),\ldots,f(v_n)\}
    \]
    is \linearlyindependent in \(W\). Indeed, if
    \[
        \lambda_1f(v_1)+\cdots+\lambda_nf(v_n)=0_W,
    \]
    then injectivity forces
    \[
        \lambda_1v_1+\cdots+\lambda_nv_n=0_V,
    \]
    so all \(\lambda_i\) vanish. Since \(W\) has dimension \(n\), the
    independent collection \(f(\mathcal{B})\) is a \basis of \(W\). Hence
    \(\ltimage f=W\), so \(f\) is \surjective.

    If \(f\) is \surjective, then \(f(\mathcal{B})\) generates \(W\). It has
    \(n\) vectors in an \(n\)-dimensional \vectorspace, so it is a \basis of
    \(W\). If \(v\in\ltkernel f\), write
    \[
        v=\lambda_1v_1+\cdots+\lambda_nv_n.
    \]
    Then
    \[
        0_W=f(v)=\lambda_1f(v_1)+\cdots+\lambda_nf(v_n).
    \]
    Since \(f(\mathcal{B})\) is a \basis of \(W\), all \(\lambda_i\) vanish,
    so \(v=0_V\). Hence \(\ltkernel f=\{0_V\}\), and \(f\) is \injective.
\end{snippetproof}

\end{document}
